\chapter*{Úvod}
\addcontentsline{toc}{chapter}{Úvod}

AI coding agenti dokáží autonomně implementovat funkcionalitu, psát
testy, spravovat verzovací historii a~komunikovat prostřednictvím
issues a~pull requestů. S~rostoucí autonomií těchto nástrojů roste
i~potřeba systematicky hodnotit kvalitu jejich práce. Současné benchmarky hodnotí izolované aspekty práce agenta,
především funkční korektnost. Holistický pohled, který by
zachytil proces, kvalitu kódu a~efektivitu společně, však chybí.
Přitom i~funkčně správné řešení může být v~praxi odmítnuto
kvůli nedostatkům v~procesu nebo udržovatelnosti. Praktik, který vidí
jen pass/fail výsledek, tak může přecenit řešení, které by ve skutečném
vývoji neprošlo code review nebo by se obtížně rozvíjelo.

Kvalitu práce agenta výrazně ovlivňují instrukce, které dostane.
V~praxi se etablují instrukční soubory jako \texttt{AGENTS.md}
nebo \texttt{CLAUDE.md}, které definují pracovní postup, konvence
a~omezení projektu. Tyto soubory se stávají standardním rozhraním
mezi člověkem a~agentem nejen v~softwarovém vývoji, ale
i~v~univerzálních agentních nástrojích. Jak instrukce systematicky
navrhovat a~jak měřit, zda agent dodržuje požadované praktiky,
zůstává otevřenou otázkou.

Tato práce navrhuje sadu metrik pokrývající proces, kvalitu kódu
a~efektivitu práce agenta. Na případové studii necháváme agenta
opakovaně implementovat systém upomínek faktur ze specifikace,
výsledky měříme navrženými metrikami a~na jejich základě
upravujeme instrukce mezi iteracemi. Ablacemi dále zkoumáme,
které složky instrukcí přispívají k~měřenému chování
a~které jsou redundantní.

% TODO(uvod-scope): Zvážit přidání odstavce "co práce JE a NENÍ" před přehled
% struktury. Riziko v posudku: oponent může přečíst tři cíle + AGENTS.md jako
% vedlejší produkt + ablace a usoudit, že je rozsah roztříštěný a hlavní claim
% nejasný. Explicitní delimitace (není benchmark agentů, není přenositelná
% instrukční sada, není srovnání modelů; JE návrh metrického postupu
% demonstrovaný na case study) tohle riziko ostříhne hned na začátku.
%
% TODO(uvod-meta): Tři úrovně práce (systém upomínek = testbed, agent = subjekt
% měření, instrukce = nezávislá proměnná) nejsou v úvodu nikde explicitně
% oddělené. Věta "necháváme agenta opakovaně implementovat systém upomínek"
% čtenáři dovoluje myslet si, že příspěvkem je buď systém upomínek, nebo agent
% samotný. Doplnit jednu větu, která tuhle meta-strukturu pojmenuje --
% klíčové pro to, aby komise pochopila, na jaké úrovni je práce přínosem.
%
% TODO(uvod-feasibility): Úvod nikde nesignalizuje, že jde o feasibility studii
% na jednom modelu a jednom projektu. Závěr to má (ř. 46--48), ale komise
% čte úvod první. Bez tohoto framingu hrozí, že čtenář očekává generalizovatelné
% výsledky a v závěru je zklamán. Doplnit krátkou větu o povaze case study
% a analytické (ne statistické) generalizaci do odstavce s cíli nebo do scope.

Práce nejprve vymezuje problém a~cíle, poté shrnuje teoretická
východiska z~oblasti softwarového inženýrství, AI agentů
a~měření kvality. Na jejich základě navrhuje metodiku, kterou
ověřuje na případové studii. Kapitola Vyhodnocení a~diskuse interpretuje
výsledky ve vztahu k~cílům a~porovnává je s~existujícím výzkumem.
